\section{Introduction}\label{sec:introduction}

Families insure themselves against job loss. When a husband loses his job, his wife can work more, and this added
worker effect is a central part of the insurance a family has before it draws on savings and public transfers
\citep{lundberg1985added,cullen2000does,stephens2002worker,shore2010better,halla2020job,birinci2019spousal,mankart2017household,bacher2025joint}.
The leading estimate is that of \citet{blundell2016consumption}: family labor supply provides 63\% of the insurance
against permanent shocks to the husband's wage, the government 20\% and savings 17\%; rerunning their code, I
reproduce their estimates exactly. That insurance needs two things the estimate takes as given: a second earner whose
job is not lost with the first, and a marriage that survives the shock. This paper shows that neither can be taken as
given. When a husband loses his job, his wife is ten times as likely to lose hers in the same month; the household is
then left with no earner, and in those months the Earned Income Tax Credit, designed to support the most vulnerable
working families, pays it nothing, because it pays only on earnings and only once a year. And a job loss of either spouse ends marriages
\citep{charles2004job}, most often when both spouses lose their jobs.

Why has this been missed? \citet{ganong2026earnings} show that most US workers' pay moves substantially from month to
month, that the movements come largely from hours set by firms rather than from wages, and that they are transitory
enough to be masked in annual data; they set aside unemployment and job changes, which, they note, would only add to
the instability. This paper studies that margin at the level of the household: the job itself is lost, and both
spouses lose theirs together. The two literatures meet here. In the two-year changes of
\citet{blundell2016consumption}, joint displacements are 0.3\% of couple-years, too rare for an estimator built on
average covariances to see, while in monthly data they occur ten times as often as independent risks imply. And if
hours and jobs are set by firms, part of what an annual framework reads as the second earner's insurance response is a
shock to the second earner: family insurance is weakest in the month of the shock.

The fact is measured to the month. In thirty years of the Current Population Survey (CPS), 1.3 million married couples in which both spouses work, a wife loses her job in
5.45\% of the months in which her husband loses his and in 0.54\% of the months in which he keeps it. The excess, 4.0
percentage points (standard error 0.002), is the same when couples are compared within the same state and month and
when wives are compared within the same state, month and industry (3.7 points), so it is not the local business cycle.
It is present in all 51 states, larger in recessions, and concentrated where spouses share a workplace: in the same
detailed industry, 23\% of wives lose their job in the month their husband does. A shared employer carries a third
to a half of the excess: couples in the same detailed industry, occupation and class of worker, the closest the public
data come to a common employer, account for about a third of it, and couples in the same industry for more than half. It is not an artefact of who reports
(3.8 points whether the wife or the husband reports the other's job loss), of temporary layoffs (2.7 points without
them) or of the season, and a joint permanent job loss---both spouses still unemployed two months later---is 13 times as
likely in the month the husband loses his job as in other months.

The fact and the estimate of \citet{blundell2016consumption} are not in conflict. Their shocks are innovations to the
hourly wages of husbands who keep working, in couples who stay together, while a job loss changes partners: the share
of displaced fathers still living with their wife falls by 4.5 percentage points within three years for Black fathers
and by 1.9 points for White fathers. Two omissions bear on the safety net: the Earned Income Tax Credit in their
schedule pays nothing in the state in which both spouses have lost their jobs, which a framework built on positive
hours does not contain, and the families who rely most on food stamps are largely outside the sample, although 71\% of
the insurance of the food-stamp-eligible households who remain comes from the government. Family labor supply insures
the wage shocks of stable, continuously employed couples; it fails against the shock at which the second earner's job
is lost with the first---which is why the joint job loss matters for the design of welfare programs.

The closest model is the joint search model of \citet{bacher2025joint}. In CPS couples in which one spouse works and the
other does not, they find that the non-working spouse is 5.6 percentage points more likely to enter the labor force
when the working spouse loses his job, 7.1 points at ages 26--35 and 1.4 points at 56--65, and they explain the decline
with age in a monthly life-cycle model of two spouses who search jointly, save and accumulate human capital, with
unemployment insurance. Their evidence concerns the spouse who can enter work; this paper concerns the spouse who is
lost, in couples in which both work. Separation shocks in their model are independent across spouses, and they note
that correlated shocks, from couples sorting into the same industries and occupations, are left out. With independent
risks, the probability that both spouses lose their jobs in the same month is close to the product of their separate
risks, about a tenth of what the CPS shows, so a model of that kind understates by construction the state in which the
household has no earner. This paper measures the correlation and builds it into a model with the structure of
\citet{bacher2025joint}: the spouses' displacement risk has a common component, larger for couples who share an employer
or an industry, and a permanent component for couples who share risk; the safety net includes the Earned Income Tax
Credit and food stamps as well as unemployment insurance; the marriage can end after a job loss; and all parameters are
estimated jointly. In the same CPS panel I reproduce their added worker effect (5.6 points) and its decline with age, from 6--7 points
before 45 to 1.2 points at 56--62; the joint job loss of spouses who both work does not decline with age (an excess of
2.8 points at 26--35 and 2.7 at 56--62 outside temporary layoffs). The spouse who can enter work is available mainly
when young; the spouse who is lost is lost at every age.

Family insurance is undiversified: a wife's job-loss hazard rises tenfold in the month her husband loses his job, the
excess is carried by shared workplaces rather than the local cycle, and because the second earner is lost exactly
when she is needed, average added-worker estimates understate what wives do and overstate what families can absorb; a
two-earner model with correlated displacement risk is built to show what this does to consumption, children and the
ranking of work-conditioned versus unconditional transfers.

The joint job loss is the paper's headline. It changes how the added-worker literature should be read and how family
insurance should be modelled. It is not itself a result about children or intergenerational mobility: children and
mobility are the stakes, which the model delivers.

The paper follows the shock in the order in which it reaches the family. Section~\ref{sec:jointloss} establishes the
fact. The CPS measures it to the month: the wife's excess risk is 0.6--0.7 points in the two months before her
husband's job loss, 4.0 points in the month of it and 0.6--0.8 points in the two months after. The Panel Study of
Income Dynamics (PSID) measures it with each spouse's own reason for the job ending. Around a husband's displacement
his wife's own displacement rises by 5.2 percentage points, from a base of 5.8\%; about 2 points of this are a
permanent trait of couples who share risk, since wives of men who will be displaced six to ten years later are already
more often displaced, and 3 points are contemporaneous. The shocks are shared at the workplace. The wife's own employer
closes in the year her husband's does, and the excess is 20 points larger when the spouses work in the same sector and
28 points larger in the same three-digit industry. It does not move with the local unemployment rate. The tetrachoric
correlation of the spouses' job losses is 0.24, and 0.50 when they share a sector. The CPS Displaced Worker
Supplement, 1984--2024, which records each adult's reason for a job loss in 417,269 couples, gives the same answer: a
wife is 4.3 times as likely to be displaced in the year her husband is, joint plant closings are almost ten times as
frequent as independent closings would imply within the same year, and the excess is the same in both races; and the
SIPP, which follows 50,692 couples month by month, gives a same-month ratio of 7.7.

Section~\ref{sec:household} asks what the household loses and what it does. A husband's displacement lowers his labor
income by 17\% (Black fathers) to 24\% (White fathers) over the following three years and family income by at least as
much. Wives work 51 more hours a year on average, 112 more when the youngest child is under six. Wives who were not
working start looking for work at once, but their employment takes about six months to build up. In the CPS the
probability that they are unemployed and looking rises by 8.8 points in the month after their husband's job loss, while
their employment rises by 2.2 points after one month and 4.2 after three; in the SIPP, which follows them for three
years, looking for work jumps at once and employment builds from 1.9 points in the month of the husband's separation to
5.4 points after six months. In the month in which the household's income falls most, its second earner is
more likely to lose her own job than to add work. This is the timing that average added-worker estimates miss: they
net the wives who work more against the wives who lose their own jobs to the same shock. The joint loss is small in
the household's average income loss (about \$260 a year, 2\% of the family's loss) but large where it bites, in the
month of the shock and among couples who share a sector, and it is absent from models of family insurance. The job
loss also strains the marriage: the share of displaced fathers still living with their wife falls by 4.5 percentage
points within three years for Black fathers and 1.9 points for White fathers. The wife's own job loss is followed by
break-up as often (4.6 points within three years, not different from the husband's), so a job loss on either side can
end the arrangement that is meant to insure it, although part of the response after her layoffs may reflect who is
laid off. Food stamp receipt rises by about four
points; AFDC/TANF barely moves.

Which part of the safety net can insure the family follows from the joint job loss. The Earned Income Tax Credit, the
largest cash transfer to families with children, pays only when someone in the household works, so it fails exactly
when both spouses lose their jobs; programs that do not require earnings, such as food stamps, are the ones that can
insure the family in that state. Section~\ref{sec:children} reports what the PSID can say about the children: children
of displaced fathers reach lower incomes and less schooling as adults, but displaced families were worse off before the
job loss, and the associations are robust only for White children's schooling. The paper uses them as bounds, not as
estimates of the effect. Section~\ref{sec:model} sets out a two-earner model with correlated displacement risk, a
search friction, savings, the marriage and the safety net, estimated on these moments, and the policy experiments that
the correlation makes informative: correlated against independent shocks, unemployment insurance for both spouses, an
earnings subsidy against an unconditional transfer of equal cost, and childcare.

The paper makes four contributions. The first is the fact: family insurance is undiversified, because the second
earner's job is lost to the same shock and in the same month as the first earner's; the excess is measured to the month
in the CPS and with each spouse's own reason for the job ending in the PSID, it is carried by shared workplaces, and it
is equal in both races. The second is what the fact implies for the added-worker literature
\citep{cullen2000does, stephens2002worker, halla2020job}: average estimates net the wives who work more against the
wives who lose their own jobs, so they understate what wives do and overstate what families can absorb in the month of
the shock. The third concerns the marriage that the insurance relies on. A job loss of either spouse raises
divorce \citep{charles2004job}, and studies of job loss follow one displaced worker at a time
\citep{eliason2012lost, athey2026heterogeneous}. Measuring both spouses' losses on one sample with one design, I
confirm over 1984--2021 that the wife's own job loss ends marriages as often as her husband's, and show that the
response is largest when both spouses lose their jobs (9.5 points fewer couples still together six years on, against
4.1 points when only he loses his). The fourth carries the facts into a
model of family insurance \citep{birinci2019spousal}, which is built to measure what undiversified insurance does to
consumption, to children and to the ranking of work-conditioned transfers against unconditional ones. Section~\ref{sec:literature} sets out what is known and what is new.

Race enters as a dimension of heterogeneity. Black men lose their jobs more often (8.0\% a year against 4.8\% for
White men in the PSID) and are the first to lose them when the economy weakens \citep{couch2010last}, and Black families
hold far less wealth \citep{derenoncourt2023wealth}, so the household's insurance matters more for them. The joint job
loss itself is the same in both races (an excess of 4.4 points for Black and 3.8 points for White couples in the CPS;
difference $p=0.309$); the separation response is larger in point estimate for Black families but not significantly so; and Black wives who were not working add less work
after their husband's job loss.

Appendix~\ref{app:other} reports evidence that bears on the model's assumptions but not on the headline: around a birth,
mothers' work does not respond to shocks to their husbands' industries (a precise null), mothers' work in a child's
first year is not associated with worse child outcomes in these data, and slack labor markets lower rather than raise
mothers' employment after a birth. Appendix~\ref{app:designs} reports the designs that did not identify what they were
meant to.
